\chapter{TAUTAN PENTING}
\label{chap:lampiran-a}

Tabel \ref{tab:tautan-penting} menyajikan tautan repositori dan \textit{deployment} yang berkaitan dengan penelitian ini.

\begin{table}[H]
  \caption{Tautan Penting}
  \label{tab:tautan-penting}
  \begin{tabularx}{\textwidth}{| l | X | X |}
    \hline
    \textbf{No} & \textbf{Tautan} & \textbf{Deskripsi} \\
    \hline
    1 & \url{https://github.com/Mipol2/TA-Node-RED} & Repositori Node-RED \\
    \hline
    2 & \url{https://github.com/Mipol2/TA-Firmware-ESP32} & Repositori \textit{Firmware} ESP32 \\
    \hline
    3 & \url{https://github.com/Mipol2/TA-Dasboard} & Repositori \textit{Dashboard} \\
    \hline
    4 & \url{https://github.com/Mipol2/TA-Data-Cleaner-Visual} & Repositori \textit{Dataset} Pengukuran Sampah serta \textit{Data Cleaner} dan \textit{Data Visualization} \\
    \hline
    5 & \url{https://dashboard-ta-mef7.vercel.app/} & \textit{Smart Waste Bin} IoT \textit{Dashboard} yang telah di-\textit{deploy} \\
    \hline
  \end{tabularx}
\end{table}